\section{Core Components and Agent Contracts}

\subsection{Agent Roles}

\begin{longtable}{p{0.18\textwidth}p{0.35\textwidth}p{0.18\textwidth}p{0.18\textwidth}}
\toprule
\textbf{Agent} & \textbf{Core responsibility} & \textbf{Inputs} & \textbf{Outputs} \\
\midrule
Orchestrator & Plans workflow, chooses agents, manages state, resolves dependencies, and aggregates final response. & User request, dataset metadata, agent results. & Task plan, agent calls, response package. \\
\addlinespace
Input Parser & Detects format, parses JSON/YAML/CSV, extracts field names, row counts, sample values, and parse errors. & Raw file or text. & Parsed object, format, parser report. \\
\addlinespace
Schema & Infers schema, compares against known schemas, identifies required and optional fields. & Parsed data, retrieved schemas. & Schema profile, candidate match, confidence. \\
\addlinespace
Validation & Checks syntax, types, missing values, duplicate keys, invalid rows, and schema violations. & Parsed data, schema profile. & Validation report, issue list, severity. \\
\addlinespace
Transformation & Executes JSON/YAML/CSV conversion, field mapping, normalization, and approved cleaning. & Validated data, plan, rules. & Converted output, diff, metadata. \\
\addlinespace
RAG Retrieval & Searches knowledge base for rules, examples, data dictionaries, and documentation. & User task, fields, schema candidates. & Retrieved context with source IDs. \\
\addlinespace
Explanation & Produces user-facing summary of the dataset, transformation, validation findings, and caveats. & Output, validation report, retrieved context. & Explanation, glossary, limitations. \\
\addlinespace
Safety and Audit & Checks prompt injection risk, sensitive data exposure, output validity, tool permissions, and logs. & Input, tool calls, outputs. & Audit log, safety flags, approval requirements. \\
\addlinespace
Clinical Explainability & Builds evidence-grounded healthcare explanations, uncertainty summaries, limitations, and clinician review prompts. & Retrieved evidence, validation trace, output, clinical safety flags. & Explanation report, claim support map, review recommendation. \\
\addlinespace
Human Review & Requests approval or clarification when actions are risky or ambiguous. & Proposed action, issue report. & Approval, edit, rejection, clarification. \\
\bottomrule
\end{longtable}

\subsection{Shared Workflow State}

All agents should operate on a shared workflow state object. This prevents inconsistent interpretations and gives the system a single auditable source of truth.

\begin{verbatim}
{
  "workflow_id": "uuid",
  "user_instruction": "Clean this CSV, convert it to JSON, and explain anomalies",
  "input_format": "csv",
  "target_format": "json",
  "dataset_ref": "storage://datasets/input-001",
  "schema_profile": {},
  "retrieved_context": [],
  "validation_report": {},
  "transformation_plan": {},
  "human_approval": null,
  "output_ref": null,
  "audit_events": []
}
\end{verbatim}

\subsection{Agent Response Contract}

Agent-to-agent communication should be structured rather than free-form. Each agent response should include:

\begin{tabularx}{\textwidth}{p{0.31\textwidth}X}
\toprule
\textbf{Field} & \textbf{Purpose} \\
\midrule
\texttt{status} & One of \texttt{success}, \texttt{warning}, \texttt{error}, or \texttt{needs\_human\_review}. \\
\texttt{summary} & One short natural language description of what the agent found or changed. \\
\texttt{data} & Structured payload, never the only place where a critical decision is hidden. \\
\texttt{confidence} & Numeric or categorical confidence, especially for schema matching and repair suggestions. \\
\texttt{issues} & Errors, warnings, policy flags, or schema problems with severity. \\
\texttt{next\_recommended\_action} & Suggested next step for the orchestrator. \\
\bottomrule
\end{tabularx}

\subsection{Workflow Example}

\begin{figure}[H]
  \centering
  \resizebox{\textwidth}{!}{\begin{tikzpicture}[
  font=\small,
  >=Latex,
  wfstep/.style={draw, rounded corners=4pt, minimum height=0.8cm, text width=2.25cm, align=center, inner sep=5pt},
  gate/.style={draw, diamond, aspect=2.0, text width=1.55cm, align=center, inner sep=2pt},
  arrow/.style={-{Latex[length=2mm]}, thick}
]

\node[wfstep, fill=sfblue!8, draw=sfblue] (submit) {User submits CSV + instruction};
\node[wfstep, fill=sfgray, draw=sfdark, right=0.55cm of submit] (parse) {Parse and profile data};
\node[wfstep, fill=sfcyan!8, draw=sfcyan, right=0.55cm of parse] (retrieve) {Retrieve schema, rules, examples};
\node[wfstep, fill=sfgray, draw=sfdark, right=0.55cm of retrieve] (validate) {Validate fields, rows, types};
\node[gate, fill=sforange!12, draw=sforange, below=0.85cm of validate] (gate) {Risky or ambiguous?};
\node[wfstep, fill=sforange!12, draw=sforange, left=0.85cm of gate] (review) {Ask human to approve, edit, or reject};
\node[wfstep, fill=sfgreen!9, draw=sfgreen, right=0.85cm of gate] (transform) {Run approved rule-based conversion};
\node[wfstep, fill=sfpurple!9, draw=sfpurple, right=0.55cm of transform] (explain) {Generate grounded explanation};
\node[wfstep, fill=sfblue!8, draw=sfblue, right=0.55cm of explain] (return) {Return output, report, audit trail};

\draw[arrow, sfblue] (submit) -- (parse);
\draw[arrow, sfdark] (parse) -- (retrieve);
\draw[arrow, sfcyan] (retrieve) -- (validate);
\draw[arrow, sfdark] (validate) -- (gate);
\draw[arrow, sforange] (gate) -- node[above] {yes} (review);
\draw[arrow, sforange] (review) -- (transform);
\draw[arrow, sfgreen] (gate) -- node[above] {no} (transform);
\draw[arrow, sfgreen] (transform) -- (explain);
\draw[arrow, sfpurple] (explain) -- (return);

\node[draw=sfred, rounded corners=3pt, fill=sfred!7, text width=2.7cm, align=center, below=0.65cm of review] (block) {Reject unsafe action, preserve original data};
\draw[arrow, sfred] (review) -- (block);
\end{tikzpicture}
}
  \caption{Human-reviewed workflow for cleaning a messy CSV, converting it to JSON, and generating a grounded explanation.}
  \label{fig:workflow}
\end{figure}
